\chapter{Introduction}\label{chap:introduction}

Equatorial Plasma Bubbles (EPBs) are a major space weather hazard triggered by the Rayleigh-Taylor instability in the post-sunset equatorial ionosphere. Because these large-scale plasma depletions cause significant amplitude and phase scintillations in trans-ionospheric radio signals, they can severely degrade Global Navigation Satellite Systems (GNSS) and satellite communications. With our growing reliance on these technologies, continuous and real-time monitoring of EPBs has never been more important \cite{zhong2025automatic,okoh2025bootstrapping,githio2024machine}.

For decades, tracking these bubbles meant relying heavily on All-Sky Imagers (ASIs) and having experts manually sift through vast amounts of high-resolution visual data \cite{okoh2025bootstrapping,zhong2025automatic}. Though this human-centric approach laid the groundwork for our current understanding, it demands immense labor and remains quite vulnerable to environmental interference, like shifting lunar illumination or passing clouds. Today, the sheer volume of observational data makes manual processing a bottleneck. Machine Learning (ML) stands out as a clear alternative to automate this task. However, feeding raw optical imagery directly into standard networks often drives up computational costs and forces models to battle background noise, highlighting the need for smarter, physics-aware detection strategies \cite{githio2024machine}.

In this work, we propose a machine learning architecture designed specifically around the physics of all-sky airglow imagery. Rather than treating the images simply as generic grids of pixels, the pipeline applies targeted preprocessing to strip away non-uniform illumination and boost local contrast, followed by a dimensionality reduction strategy that preserves the natural two-dimensional structure of the data. An ensemble classifier then evaluates the extracted features, merging the predictive strengths of several base algorithms to identify plasma structures reliably.

Ultimately, our goal is to validate a streamlined detection system for EPBs that respects the underlying physics of the data while remaining interpretable, fast, and computationally light enough for real-time monitoring.

The remainder of this document is organized as follows. Chapter~\ref{chap:objectives} presents the general and specific objectives that guide this research. Chapter~\ref{chap:methodology} describes the full technical methodology, covering dataset acquisition, the physics-informed preprocessing pipeline, feature extraction via 2DPCA, the stacking ensemble classifier, and the evaluation strategy. Chapter~\ref{chap:results} reports the preliminary results obtained so far, including cross-validation performance, model interpretability analysis, and temporal stability tests. Chapter~\ref{chap:schedule} outlines the execution schedule for the remaining phases of the dissertation. Finally, Chapter~\ref{chap:conclusion} summarizes the main findings and discusses directions for future work.